% Source: https://artofproblemsolving.com/community/c3h2126371_prime_numbers
% Stable identity: aops-topic-2126371-q01
The statement is \textbf{false}. For the first primes,
\[2^2-1=3,\qquad 2^3-1=7,\qquad 2^5-1=31,\qquad 2^7-1=127\]
are all prime ($127$ is not divisible by $2,3,5,7$ or $11$, and $13^2>127$). But for $p=11$,
\[2^{11}-1=2047=23\cdot89 .\]
So the smallest counterexample is $\boxed{p=11}$.

(The factor $23$ is natural to try: $2^{11}=2048=23\cdot89+1$. Another quick check is that $2^{11}\equiv1\pmod{23}$ because $2^{11}=\left(2^{5}\right)^2\cdot2=32^2\cdot2\equiv9^2\cdot2=162\equiv1\pmod{23}$.)

\emph{Remark.} The converse statement is true: if $2^n-1$ is prime then $n$ is prime, because for $n=ab$ with $1<a<n$ the number $2^a-1$ is a proper divisor of $2^n-1$. Primes of the form $2^p-1$ are called Mersenne primes.
